\subsection{Notes to self}

\Joris{
\begin{enumerate}
  \item A new idea: Leihao suggested to mimic input nodes by treating them as output nodes fully connected by UNdirected edges.
    It could be that with that convention, the extended FCI rules become identical to the FCI ones, except for adding background knowledge.
    We could try adding the background knowledge by imposing dummy undirected four-cycles on the input nodes.
    The only issue could be that rules 5 and 6, which lead to the implementation of the background knowledge, are too late in the game.
    So: if R1-R5 do not care about the background knowledge, then it appears we have reduced extended FCI to vanilla FCI with dummy nodes!
    Hm, at first sight rules R0,R1,R2 do not care about the background knowledge, but R3 might (before background knowledge is applied,
    $l$ might be in $J$, but afterwards, $l$ cannot be in $J$ because of the circles; however, this is ``the wrong way around''; R3
    would fire more often if we postpone background knowledge, so actually doing that may just make the algorithm MORE complete!).
    \begin{center}\begin{tikzpicture}
    \begin{scope}
      \node at (-4,0) {(a)};
      \node[ndout] (Xi) at (-1.5,0) {$i$};
      \node[ndout] (Xj) at (0,-1.5) {$j$};
      \node[ndout] (Xk) at (1.5,0) {$k$};
      \node[ndout] (Xl) at (0,1.5) {$l$};
      \draw[suh] (Xi) edge (Xj);
      \draw[suh] (Xk) edge (Xj);
      \draw[suo] (Xi) edge (Xl);
      \draw[suo] (Xk) edge (Xl);
      \draw[suo] (Xl) edge (Xj);
    \end{scope}
    \end{tikzpicture}\end{center}
    Same holds for R4: for this to fire, it expects a circle at $j$ (which might be an input node), so postponing background knowledge
    may make this fire more often!
    (NOTE: this analysis also suggests more incompleteness cases of the current extended FCI...
    Indeed, R3 seems to miss a case where as it assumes that the ONLY definite non-colliders are unoriented unshielded triples;
    but also nodes in $J$ are definite non-colliders but with background knowledge applied they may not be part of UNORIENTED unshielded triples.
    This appears to make the current FCI algorithm incomplete!
    We should add an OR clause to the antecedent of R3: ``$i,l,k$ is unshielded triple or $l \in J$''.
    Or perhaps postpone background knowledge application until after R3.
    Something similar may hold for other rules as well that look at unshielded triples!)
    R0 is clearly invariant under applying-background-knowledge-before-or-after.
    R1 as well.
    In general, note that arrowheads are invariant.
    So in R2, only $i$ could be in $J$. One then easily sees that R2 is invariant.
    For R3, $i,k,l$ could be in $J$. R3 doesn't care about the edge marks near $i$ or $k$. With the above extension of R3 (so it catches the case $l \in J$ with background knowledge applied i.e.\ $i \sut l \tus k$ instead of $i \suo l \ous k$) it becomes invariant too.
    For R4, only $j$ or $k$ could be in $J$. The current formulation may also forget to orient $i \hus j$ if $j \in J$ and background knowledge has been applied to $j$ i.e.\ $i \sut j$ is in the antecedent;
    however, if background knowledge is applied before R1 then R1 will catch this. So current FCI version is complete enough in that respect.
    \textbf{So it appears this strategy may just work and make extended FCI (with additional extension of R3) complete!!!}
    To check whether this all works out requires a bit of attention.
    Perhaps the safest strategy is to first come up with a formal correspondence between our current $\sigma$-PAGs with input nodes,
    and normal PAGs with input nodes replaced by fully-connected-by-undirected-edges output nodes.
    Are both representations equivalent?

    Interesting: the rules that I and Tom thought needed to be added (point \ref{fci_notes:additional_rules} below) are redundant,
    at least if we exclude cycles.
    For $i \tuo j \hus k$ with $i \in J$, orient $i \tuh j$: there exists no valid MAG with $i \tut j \hus k$, so by arrowhead
    completeness of FCI, this will get oriented.
    For $i \tuo j \ous k$ with $i \in J$, $k \notin J$, $i,k$ non-adjacent, orient $j \tus k$.
    If there would exist a valid MAG with $j \hus k$, then it must have $i \tuh j \hus k$, but that contradicts that rule $\mathcal{R}_0$
    has been exhaustively applied. So by tail completeness of FCI, this will get oriented.

    Suppose we have run vanilla FCI with dummy variables augmented for implementing the background knowledge.
    By its completeness, if there remains $i \tuo j$ with $i \in J$, then apparently there is a MAG with $i \tut j$ and one with $i \tuh j$.
    In the first case, we may still try to exploit that $j \notin \Anc(i)$. 
    So we have to check all FCI rules which expects an arrowhead $i \tuh j$ whether it only relies on $j \notin \Anc(i)$ (because we do know $j \in \Anc(S)$).
    In the second case, there is nothing left to exploit.

  \item Perhaps we can do all of this without exogenous input nodes by instead looking at $P_{M}(X_O \mid X_S \in \xi_S, X_J)$.
Then we can treat the $J$ nodes as output variables, and still make use of qualitatively similar causal background knowledge
on the relationship between $J$ nodes and the other nodes. For example, we might get undirected edges $j_1 \tut j_2$ in the PAG
if $j_1$ and $j_2$ have a common descendant in $S$. We can still try to deduce whatever we can deduce from the conditional
independence model. The difference between boxes and circles then just becomes the difference between $\forall x_J$ and 
$P(X_J)$-a.s. Idea: study RCT in the presence of selection bias.

\item An alternative setup could be to assume the existence of a simple SCM $M = \lt \Sin, \Send, \Srnd, \CX, P, f \rt$, and to
assume that the exogenous input variables $J$ and a subset of endogenous and exogenous random variables $O \subseteq V \cup W$ are observed. 
Note: for the observed exogenous random variables, we do have background knowledge, namely that they are independent and exogenous. 
If we have that background knowledge, why don't we use it? 
So it seems more sensible to assume that only a subset of the endogenous variables, and the exogonous input variables, are observed. 
This suggests to take $V^+ = V \dcup S \dcup L$ where only the ones in $V$ are observed.

 \item At some point I thought one would wish to demand that $i \notin J$ or $j \notin J$ for the existence of a $\sigma$-inducing path between $i$ and $j$ given $S$, but it actually appears more natural to have $\sigma$-inducing walks/paths between any pair of nodes in $J$.

 \item In Proposition~\ref{prp:sigma-inducing_path_properties}, we could add a fifth equivalence:
    $i \nisPerp_G j \given S \cup (J \sm \{j\}) \cup Z$ for all $Z \subseteq V \sm \{i,j\}$.
   This might be true, but it seems we don't need it.

  \item If we condition on a common child of two exogenous input nodes, the two input nodes may no longer be variationally independent. It seems we only need that assumption for a \emph{causal} interpretation of the input nodes (as `independent' causes). So if we allow for selection bias, then we already loose the ability to interpret an ancestral relation from input node to output nodes as causal. So I guess we should indeed be quite careful when combining selection bias and input nodes!!

  \item It seems that a rather neat solution for the (non)causal interpretation (for both input nodes and output nodes) is obtained as follows.
    We explicitly mark each nodes as `causal' / `not-necessarily-causal'. 
    In the FCI PAG, each node that has an arrowhead is marked as `causal'.
    All other nodes are marked as `not-necessarily-causal'.
    This is complete in the sense that all nodes that are an identifiable non-ancestor of $S$ are then marked `causal', 
    and all nodes that are not identifiable non-ancestors of $S$ are then marked `not-necessarily-causal'.
    Input nodes will never get arrowheads, so for these we could make use of background knowledge to mark some of them as `causal' (by default they would all be `not-necessarily-causal').
    (A hacky way to make FCI use this background knowledge would be to add a dummy node and edge $j \huo j'$ to an input node $j$).

  \item We now dealt with interactions between selection bias and input nodes by reformulating several of the FCI rules to explicitly exclude certain cases.
    Another option might be to just take all $J$ nodes to be adjacent. In both cases, we still need to investigate which rules may be added to get a more complete algorithm.
    We consider several exceptional cases (see also Figure~\ref{fig:fci_input_nodes_exceptions}).
    \begin{enumerate}
      \item[$\C{R}0$] (``if $i \sus j \sus k$ in $H$ with $i \notin J$, and $i$ and $k$ are not adjacent in $H$, then orient $i \suh j \hus k$ if $j \notin \sepset(\{i,k\})$'')
        This rule is invalid for $i,k \in J$, see Figure~\ref{fig:fci_input_nodes_exceptions}(a). Additionally, the question is what $\sepset(\{i,k\})$ would be in that case. 
        Perhaps it is more natural after all to consider the nodes in $J$ as connected...
      \item[$\C{R}1$] (``if $i \suo j \hus k$ in $H$ with $i \notin J$, and $i$ and $k$ are not adjacent in $H$, then orient $i \hut j$'')
        This rule is invalid for $i,k \in J$, see Figure~\ref{fig:fci_input_nodes_exceptions}(b).
      \item[$\C{R}3$] (``if $i \suh j \hus k$ and $i \suo l \ous k$ and $l \suo j$ in $H$ with $i \notin J$, and $i$ and $k$ are not adjacent in $H$, then orient $l \suh j$'') This rule is invalid if $i,k \in J$, see Figure~\ref{fig:fci_input_nodes_exceptions}(c).
      \item[$\C{R}7$] (``if $i \suo j \out k$ in $H$ with $i \notin J$, and $i$ and $k$ are not adjacent in $H$, then orient $i \sut j$'')
        This rule is invalid for $i,k \in J$, see Figure~\ref{fig:fci_input_nodes_exceptions}(d).
    \end{enumerate}
    There are two ways to fix these rules: treating all $J$ nodes as adjacent, or (the current way) treat all $J$ nodes as non-adjacent but explicitly exclude the case $i \in J$.
    The former has the additional advantage that then the existence of a $\sigma$-inducing walk between two nodes (even two input nodes) becomes again equivalent to adjacency in the skeleton.
    \begin{figure}
      \begin{tikzpicture}
        \begin{scope}
          \node at (0,1) {$\C{R}0$:};
          \node[ndint] (J1) at (-1,0) {$i$};
          \node[ndint] (J2) at (1,0) {$k$};
          \node[ndout] (X) at (0,-1.5) {$j$};
          \node[ndout] (S) at (0,-3) {$s$};
          \draw[tuh] (J1) edge (X);
          \draw[tuh] (J2) edge (X);
          \draw[tuh] (X) edge (S);
        \end{scope}
        \begin{scope}[xshift=4cm]
          \node at (0,1) {$\C{R}1$:};
          \node[ndint] (I) at (-1,0) {$i$};
          \node[ndint] (K) at (1,0) {$k$};
          \node[ndout] (J) at (1,-1.5) {$j$};
          \node[ndout] (L) at (-1,-1.5) {$l$};
          \draw[tuh] (I) edge (J);
          \draw[tuh] (K) edge (J);
          \draw[tuh] (L) edge (J);
          \draw[tuh] (I) edge (L);
        \end{scope}
        \begin{scope}[xshift=8cm]
          \node at (0,1) {$\C{R}3$:};
          \node[ndint] (I) at (-1,0) {$i$};
          \node[ndint] (K) at (1,0) {$k$};
          \node[ndout] (J) at (-1,-1.5) {$j$};
          \node[ndout] (L) at (1,-1.5) {$l$};
          \draw[tuh] (I) edge (J);
          \draw[tuh] (K) edge (J);
          \draw[tuh] (K) edge (L);
          \draw[tuh] (I) edge (L);
          \draw[tuh] (J) edge (L);
        \end{scope}
        \begin{scope}[xshift=12cm]
          \node at (0,1) {$\C{R}7$:};
          \node[ndint] (I) at (-1,0) {$i$};
          \node[ndint] (K) at (1,0) {$k$};
          \node[ndout] (J) at (-1,-1.5) {$j$};
          \draw[tuh] (I) edge (J);
          \draw[tuh] (K) edge (J);
        \end{scope}
      \end{tikzpicture}
      \caption{Graphs $G$ as counterexamples to the vanilla FCI rules when treating input nodes as non-adjacent.\label{fig:fci_input_nodes_exceptions}}
    \end{figure}

  \item Some additional orientation rules to exploit the background knowledge:\label{fci_notes:additional_rules}
    \begin{itemize}
      \item if $i \tus j \hus k$ in $H$ with $i \in J$, orient $i \tuh j$
      \item if $i \tus j \ous k$ in $H$ with $i \in J$ and $i, k$ not adjacent (update: and $k \notin J$?!), orient $j \tus k$ (this one was suggested by Tom)
          (If we additionally somehow know that $j \notin \Anc_G(S)$ then we could orient $i \tuh j \tuh k$; this looks like a combination of the previous item and R1).
          (If instead we additionally somehow know that $j \in \Anc_G(S)$ then we could orient $i \tut j \tus k$. This looks like a possible extension of $\C{R}6$, but we should investigate all the possibilities of inferring $j \in \Anc_G(S)$. It seems that at the stage where $\C{R}6$ is executed, the only possible way to know that $j \in \Anc_G(S)$ is if it is introduced by $\C{R}5$, which requires the presence of an edge $l \tut j$. Another possibility would be additional background knowledge on $j$. FInally, if we were to add this rule as $\C{R}6b$, then also this one may have oriented an edge $i' \tut j$ for some $i' \in J$.)
        \item Note that all this implies that if $j_1 \sus v \sus j_2$ with $j_1,j_2 \in J$ (and $v \in V$) then 
          we can orient either $j_1 \tuh v \hut j_2$ or $j_1 \tut v \tut j_2$, but not in any other way.
          In the latter case, $j_1,j_2,v$ are in $\Anc_G(S)$; in the former case, none of them is in $\Anc_G(S)$.
    \end{itemize}
    Would that be all? (It looks as if these rules might correspond to the ``missing'' cases of rules 0,1,3,7?!)

    Tom suggested the following strategy:
    \begin{itemize}
      \item first run vanilla fci, using only restricted CI tests (i.e.\ the input nodes stay adjacent)
      \item apply background knowledge (including making input nodes separated)
      \item propagate background knowledge
    \end{itemize}
    For an oracle, it may be easier to get this one complete.
    Note: rules 1-5 don't overwrite the background knowledge orientation ($j \tus v$ ($j\in J, v\in V$).
    Also, rules 0 and 1 do not exploit the background knowledge orientation.
    Rule 2 does seem to be able to exploit the background knowledge orientation in certain cases.
    However, by adding the rule 
      \begin{center}``if $j \sus i \hus k$ in $H$ with $j \in J$, orient $j \tuh i$''\end{center}
    this case will also be caught (so this suggests to add this as $\C{R}2$b!).
    Rules 3,4,5 don't need tails so also these commute with the background knowledge orientation.
    Now rule 6 is interesting: it seems as if this can be used to incorporate the background knowledge orientation by adding dummy nodes
    (indeed, for each $j \in J$ add a $j'$ and connect it as $j' \tut j$).
    Rule 6 itself also commutes with the background knowledge orientation step.
    Rule 7 also commutes with the background knowledge orientation step.
    Also rule 8. Also rule 9. Also rule 10.
    This is great news: it suggests that by adding this rule 2b, we can apply the background knowledge at any arbitrary point in the algorithm.
    By introducing dummy variables as above, rule 6 already catches the background knowledge orientations, hence they become redundant.
    So after adding rule 2b, the algorithm will give the same output if we either apply the background knowledge orientation 
    explicitly without using dummy variables at any point, or if we don't but add dummy variables.
    The simplest idea would be to have dummy variables as $j' \tut j'' \tut j''' \tut j \tut j'$ for each $j \in J$ (that is, a square with $j$ being one of the edges).
    Left to check: the dummy variables shouldn't have any side effects downstream!
    Hm, there is one concern: in this way, one could even conclude that $j$ is ancestor of $S$, which is not true.
    Interesting, it feels as if we're getting closer to completeness!

  \item Hold on: it appears as if the failure of R0 trickles down to several other rules.
    When does R0 fail? In the example in Figure~\ref{fig:fci_input_nodes_exceptions}, also $i$ and $k$
    are ancestor of $S$. Therefore, they should get merged! But then it is no longer a failure of R0!
    It would also suffice if we mark them as non-causal. But in that case, we would also have to 
    drop the assumption of variation independence of the input nodes.
    This all seems to be an intricate interplay.
    Note, though, that we only need to merge these input nodes if some of their joint states are
    rejected by the selection mechanism. But in order to have consistent CI tests, we already have
    to assume that each input state occurs in the data.
    So I guess that both for proving consistency, and for obtaining a causal interpretation of the
    input nodes, we have to assume variation independence of the input variables.
    However: note that for a non-causal interpretation of the input nodes, we might weaken the
    null hypothesis to consider only those values of the input variables that are realized in the data.
    In other words, we would weaken the ``for all $x_J$'' quantifier to the ``for $\mu$-a.a.\ $x_J$''
    quantifier for some measure $\mu$ (actually, we only need the null sets of the measure $\mu$).
    This is in line with what we'd get if we would condition on $x_J$: a kernel defined up to null sets on $\CX_J$.
    This weakens the requirements on the data to get consistency.
    So, it seems we end up with two types of input nodes:
    \begin{itemize}
      \item causal input nodes (not dashed), which represent an intervention, and therefore we 
        take all values in a Cartesian product
      \item non-causal input nodes (dashed), which can represent conditioning, and may lead to
        a kernel that is defined up to certain null sets
    \end{itemize}
    A simple workaround could be to \emph{assume} that the selection mechanism doesn't rule out
    any joint value of the input nodes: then we might still have a causal interpretation for the
    input nodes, and consistency of CI tests.

  \item Another thought: the current choice of non-adjacent input nodes works well for representing
    CI information, and part of the causal information. But it doesn't seem to be informative enough
    to represent the causal information between the input nodes. This was obviously precisely what
    we were aiming for: by conditioning on them, we would need to encode \emph{less} causal information.

  \item We should double-check whether in the code we are correctly using the (a)symmetric sepset!

  \item Also see file \texttt{old/fci-extensions.tex} in \texttt{jci-paper} repository.

  \item Note that in FCI we model conditioning of the type $\given X_S = x_s$ while the conditioning SCM operation considers the type $\given X_S \in \xi_S$ for some measurable $\xi_S \subseteq \CX_S$! 
    The difference is that for the former, non-collider nodes in $S$ are blocked, whereas for the latter, these are not blocked in general.
    Consider CDMG $G$: $A \tuh S \tuh B$. If we condition on $\given X_S \in \xi_S \subseteq \CX_S$ then in general there will be a dependency between $A$ and $B$, and hence this should count as a $\sigma$-inducing path as well! However, if we condition on $\given X_S=x_s$, then it will be blocked, and hence not $\sigma$-inducing.
    Would it be possible to derive an FCI version that allows conditioning on a set of values...?
    Wait! We can always add endogenous variables $T_i$ with structural equations $T_i = \ind_{\xi_i}(S_i)$. 
    Then we have that 
      $$P_M(X_V \mid X_S \in \xi_S, \doit(X_J)) = P_M(X_V \mid X_S = 1, \doit(X_J))$$
    where we marginalized out the nodes in $S$.
    For all nodes in $v \in V$, we have $v \in \Anc(T) \iff v \in \Anc(S)$.
    In this way, we reduce selection on $X_S \in \xi_S$ to selection on $X_T = 1$.

  \item There seems to be a conflict in the definition of $\sigma$-inducing path. Note that even 
    though there may or may not be a $\sigma$-inducing path between two input nodes, 
    This is why Proposition~\ref{prp:sigma-inducing_path_properties} does not consider the case $i,j \in J$.
    This mismatch can occur if input nodes are ancestors of selection variables (but perhaps also more generally? at least it requires $S \ne \emptyset$).
    SO on the one hand, we insist on giving the inpude nodes a causal interpretation. 
    However, in the presence of selection bias, we can (so far) only give the \emph{non-ancestors} of selection nodes a causal interpretation. 
    Shouldn't we therefore assume the background knowledge that each input node is not ancestor of a selection variable? 
    But, wouldn't that be a very modeling strong assumption that we would perhaps not like to make in practice?
    So this setup seems to lead us automatically to ``non-causal'' input nodes! In other words:
    even if we introduce input nodes only to represent interventions, after latent selection on a descendant the input node no longer (?) represents an intervention (!).
    For the time being, a good workaround might be to note that an input node \emph{in a $\sigma$-PAG} not necessarily has a causal interpretation, even though it does in the CDMG!
\end{enumerate}
}
